% ============================================================================
%  Problems: การนับเบื้องต้น (Counting Principles)
%  Ordered from easy (basic) → hard (POSN level).
%  Shared by exercise.tex and answer-key.tex
% ============================================================================

% --- Part 1: พื้นฐาน (Basic) ---
\examsection{ตอนที่ 1: พื้นฐาน — หลักการบวกและหลักการคูณ}

\problem{
  โรงอาหารมีข้าว 3 อย่าง และกับข้าว 5 อย่าง
  ถ้าต้องการเลือกข้าว 1 อย่าง \textbf{หรือ} กับข้าว 1 อย่าง จะมีทั้งหมดกี่วิธี?
}{
  $3 + 5 = 8$ วิธี
}

\problem{
  มีเสื้อ 4 ตัว กางเกง 3 ตัว และรองเท้า 2 คู่
  จะแต่งตัวด้วยเสื้อ 1 ตัว กางเกง 1 ตัว และรองเท้า 1 คู่
  ได้ทั้งหมดกี่แบบ?
}{
  $4 \times 3 \times 2 = 24$ แบบ
}

\problem{
  โยนเหรียญบาท 1 เหรียญ 4 ครั้ง
  จะมีผลลัพธ์ที่เป็นไปได้ทั้งหมดกี่แบบ?
  (Hint: แต่ละครั้งออกได้ 2 แบบ)
}{
  $2^4 = 16$ แบบ
}

\problem{
  ทอยลูกเต๋า 1 ลูก 2 ครั้ง
  \begin{enumerate}
    \item[(ก)] จะมีผลลัพธ์ทั้งหมดกี่แบบ?
    \item[(ข)] จะมีกี่แบบที่ผลรวมของแต้มเป็นเลขคี่?
  \end{enumerate}
  (Hint สำหรับ (ข): ผลรวมเป็นคี่เกิดจาก คู่+คี่ หรือ คี่+คู่)
}{
  (ก) $6 \times 6 = 36$ แบบ

  (ข) คู่+คี่: $3 \times 3 = 9$, \quad คี่+คู่: $3 \times 3 = 9$,
  รวม $9 + 9 = 18$ แบบ
}

\problem{
  หมายเลขโทรศัพท์ประกอบด้วยตัวเลข 0--9 จำนวน 7 หลัก
  จงหาจำนวนหมายเลขโทรศัพท์ที่เป็นไปได้ทั้งหมด
  (ตัวเลขสามารถซ้ำกันได้ และหลักแรกเป็น 0 ได้)
}{
  $10^7 = 10{,}000{,}000$ หมายเลข
}

\problem{
  ในการโยนเหรียญ 2 เหรียญพร้อมกัน 1 ครั้ง
  \begin{enumerate}
    \item[(ก)] จะมีผลลัพธ์ที่เป็นไปได้ทั้งหมดกี่แบบ?
    \item[(ข)] จงเขียนผลลัพธ์ทั้งหมดที่ได้หัวอย่างน้อย 1 เหรียญ
  \end{enumerate}
}{
  (ก) $2 \times 2 = 4$ แบบ ได้แก่ HH, HT, TH, TT

  (ข) HH, HT, TH (3 แบบ)
}

% --- Part 2: การประยุกต์ (Intermediate) ---
\examsection{ตอนที่ 2: การประยุกต์ — การเรียงสับเปลี่ยนและการจัดหมู่}

\problem{
  จงหาค่าต่อไปนี้
  \begin{enumerate}
    \item[(ก)] $4!$
    \item[(ข)] $P(7, 3)$
    \item[(ค)] $C(7, 3)$
  \end{enumerate}
}{
  (ก) $4! = 24$

  (ข) $P(7, 3) = 7 \times 6 \times 5 = 210$

  (ค) $C(7, 3) = \frac{7 \times 6 \times 5}{3 \times 2 \times 1} = 35$
}

\problem{
  จัดตัวอักษรจากคำว่า ``MATH'' มาเรียงเป็นคำใหม่
  (ไม่จำเป็นต้องมีความหมาย) จะได้ทั้งหมดกี่คำ?
}{
  $4! = 24$ คำ
}

\problem{
  มีนักเรียน 8 คน ต้องการเลือกคณะกรรมการ 3 ตำแหน่ง
  (ประธาน, รองประธาน, เลขานุการ) จะมีวิธีการเลือกทั้งหมดกี่วิธี?
}{
  $P(8, 3) = 8 \times 7 \times 6 = 336$ วิธี
}

\problem{
  มีนักเรียน 8 คน ต้องการเลือกตัวแทนไปแข่งขัน 3 คน
  (ทุกคนมีฐานะเป็นตัวแทนเท่ากัน) จะมีวิธีการเลือกทั้งหมดกี่วิธี?
}{
  $C(8, 3) = \frac{8 \times 7 \times 6}{3 \times 2 \times 1} = 56$ วิธี
}

\problem{
  มีดอกกุหลาบ 4 สี ที่แตกต่างกัน
  (แดง, ขาว, ชมพู, เหลือง) ถ้าต้องการเลือกดอกไม้ 2 สี
  (ลำดับของสีไม่สำคัญ) จะมีวิธีการเลือกทั้งหมดกี่วิธี?
}{
  $C(4, 2) = 6$ วิธี
}

\problem{
  ในการเลือกไพ่ 3 ใบจากสำรับ 52 ใบ
  โดยไม่คำนึงถึงลำดับ จะมีทั้งหมดกี่วิธี?
}{
  $C(52, 3) = \frac{52 \times 51 \times 50}{3 \times 2 \times 1} = 22{,}100$ วิธี
}

\problem{
  จัดนักเรียน 6 คนนั่งรอบโต๊ะกลม
  จะนั่งได้ทั้งหมดกี่วิธี?
}{
  $(6-1)! = 5! = 120$ วิธี
}

\problem{
  นำลูกปัด 5 สีที่แตกต่างกันร้อยเป็นสร้อยคอ
  จะได้สร้อยคอที่แตกต่างกันทั้งหมดกี่แบบ?
  (การพลิกด้านสร้อยคอถือว่าเป็นแบบเดียวกัน)
}{
  $\dfrac{(5-1)!}{2} = \dfrac{24}{2} = 12$ แบบ
}

% --- Part 3: ระดับ สอวน. (POSN Level) ---
\examsection{ตอนที่ 3: ระดับข้อสอบ สอวน.}

\problem{
  ในงานเลี้ยงมีคู่สามีภรรยา 4 คู่
  ถ้าต้องการเลือกคน 3 คนจากทั้งหมด
  \begin{enumerate}
    \item[(ก)] จะมีวิธีการเลือกทั้งหมดกี่วิธี?
    \item[(ข)] จะมีกี่วิธีที่เลือกได้สามีอย่างน้อย 1 คน?
  \end{enumerate}
}{
  (ก) $C(8, 3) = 56$ วิธี

  (ข) วิธีทั้งหมด $-$ วิธีที่ไม่ได้สามีเลย
  $= C(8,3) - C(4,3) = 56 - 4 = 52$ วิธี
}

\problem{
  นักเรียน 5 คน ได้แก่ a, b, c, d และ e จะจัดแถวเรียงเดี่ยวได้กี่วิธี ถ้า a และ b ต้องอยู่ติดกันเสมอ
}{
  48 วิธี
}

\problem{
  มีนักเรียน 8 คน เป็นชาย 3 คนหญิง 5 คน ต้องการเลือกตัวแทนไปแข่งขัน 5 คน
  (ทุกคนมีฐานะเป็นตัวแทนเท่ากัน) จะมีวิธีการเลือกทั้งหมดกี่วิธี 
  \begin{enumerate}
    \item[(ก)] ถ้าต้องมีผู้ชาย 2 คนเป็นตัวแทน
    \item[(ข)] ถ้าต้องมีหญิงอย่างน้อย 3 คนเป็นตัวแทน
  \end{enumerate}
}{
  (ก) $30$ วิธี

  (ข) $30 + 15 + 1 = 46$ วิธี
}

\problem{
  มีลูกบอล 10 ลูก ใส่ในกล่อง 3 ใบ
  โดยหลักการรังนกพิราบ จงแสดงว่าจะต้องมีอย่างน้อย 1 กล่อง
  ที่มีลูกบอลอย่างน้อยกี่ลูก?
}{
  เฉลี่ย: $\lceil \frac{10}{3} \rceil = 4$ ลูก

  ดังนั้นต้องมีอย่างน้อย 1 กล่องที่มีลูกบอลอย่างน้อย 4 ลูก
}

\problem{
  (ข้อสอบจริงปี 68 ข้อ 20) กำหนดเลขลอตเตอรี่ประกอบด้วยเลขโดด 0-9 จำนวน 4 หลัก จะได้รางวัลที่ 1 เมื่อถูกทั้ง 4 หลัก และรางวัลที่ 2 เมื่อถูก 3 หลัก จำนวนวิธีที่จะได้รางวัลที่ 2 เท่ากับเท่าใด
}{
  36 วิธี
}

\problem{
  (ข้อสอบจริงปี 67 ข้อ 32) ในวงกลมมีจุด 8 จุด ไม่ทับซ้อนกัน ลากเส้นคอร์ดเชื่อมจุดเหล่านี้ จะลากได้ทั้งหมดกี่เส้น
}{
  28 เส้น
}

\problem{
  (ข้อสอบจริงปี 67 ข้อ 33) จำนวนเต็มบวก 4 หลัก ซึ่งเป็นจำนวนคู่และไม่มีเลขซ้ำกันมีทั้งหมดกี่จำนวน
}{
  2296 จำนวน
}
